\documentclass[11pt,a4paper]{article}

\usepackage[utf8]{inputenc}
\usepackage[T1]{fontenc}
\usepackage{lmodern}
\usepackage[a4paper,margin=1in]{geometry}
\usepackage{setspace}
\usepackage{amsmath,amssymb,amsfonts,mathtools}
\usepackage{bm}
\usepackage{booktabs}
\usepackage{array}
\usepackage{multirow}
\usepackage{enumitem}
\usepackage{xcolor}
\usepackage{hyperref}
\usepackage{graphicx}

\hypersetup{
    colorlinks=true,
    linkcolor=blue!60!black,
    citecolor=blue!60!black,
    urlcolor=blue!60!black,
    pdftitle={TAME, SEG, and ChebNet Foundation Pretraining: Architecture Report},
    pdfauthor={MOLFusion2 Project}
}

\setstretch{1.12}
\setlength{\parskip}{0.4em}
\setlength{\parindent}{0pt}

\title{\textbf{A Technical Architecture Report for Multi-Modal Molecular Learning}\\
\large TAME as the Primary Model, SEG as the Bi-Modal Baseline, and ChebNet Foundation Pretraining}
\author{MOLFusion2 Research Stack}
\date{\today}

\begin{document}
\maketitle

\begin{abstract}
This report documents the architecture and training methodology of a multi-model molecular prediction stack centered on TAME (Multi-Modal Fusion with Descriptor-augmented element-wise Mixture-of-Experts), with SEG (Structured Embedding Fusion) as a strong bi-modal baseline and ChebNet foundation pretraining as the transfer backbone. We present the exact modeling assumptions, mathematical formulations, masking strategies, objective functions, and implementation-level design choices used in the codebase. The report is intended as a high-rigor technical reference suitable for scientific reporting, reproducibility, and future extension to broader benchmark suites.
\end{abstract}

\tableofcontents

\section{Scope and Contributions}
This report covers three tightly coupled components:
\begin{enumerate}[label=\arabic*.]
    \item \textbf{ChebNet foundation pretraining}: two-stage graph representation pretraining with node-level masking and graph-level descriptor prediction.
    \item \textbf{SEG}: graph+text fusion with modular fusion operators for supervised downstream tasks.
    \item \textbf{TAME}: graph+text+descriptor fusion with element-wise tri-modal routing, currently the largest and most expressive architecture in the stack.
\end{enumerate}

The primary contribution of this architecture family is the integration of:
\begin{itemize}
    \item spectral graph encoding (Chebyshev filtering),
    \item language-informed molecular context (text embeddings),
    \item classical chemistry descriptors (RDKit/Mordred style features),
    \item modality-adaptive fusion at hidden-dimension granularity.
\end{itemize}

\section{Notation and Problem Setup}
Let a molecule be represented by a graph
\begin{equation}
G = (V, E, \mathbf{X}, \mathbf{A}),
\end{equation}
where $V$ is the atom set, $E$ is the bond set, $\mathbf{X} \in \mathbb{R}^{|V| \times d_x}$ are node features, and $\mathbf{A}$ is adjacency information.

For each molecule we may have:
\begin{itemize}
    \item a text embedding $\mathbf{t} \in \mathbb{R}^{d_t}$,
    \item a descriptor vector $\mathbf{d} \in \mathbb{R}^{d_d}$,
    \item target(s) $\mathbf{y}$ (single-task or multi-task; classification or regression).
\end{itemize}

Tasks addressed in practice include:
\begin{itemize}
    \item \textbf{single-task classification} (e.g., BACE),
    \item \textbf{multi-label classification} with missing labels and masking (e.g., Tox21),
    \item \textbf{multi-task regression} (e.g., QM8).
\end{itemize}

\section{Shared Graph Backbone: ChebNet Encoder}
All three systems rely on a Chebyshev polynomial graph filter backbone. Using the scaled Laplacian $\tilde{\mathbf{L}}$, a Cheb convolution can be written as:
\begin{equation}
\mathbf{H}^{(\ell+1)} = \sum_{k=0}^{K} \mathbf{T}_k(\tilde{\mathbf{L}})\,\mathbf{H}^{(\ell)}\,\mathbf{\Theta}_k^{(\ell)},
\end{equation}
where $\mathbf{T}_k$ is the $k$-th Chebyshev polynomial, $K$ is filter order, and $\mathbf{\Theta}_k^{(\ell)}$ are learnable weights.

The encoder stacks multiple such layers and applies regularization (dropout) as configured. Graph-level pooling is then applied using a selectable operator:
\begin{itemize}
    \item mean/sum pooling,
    \item set2set pooling,
    \item attention pooling (where configured).
\end{itemize}

\section{ChebNet Foundation Pretraining}
\subsection{Rationale}
Pretraining is used to initialize chemically meaningful graph representations prior to downstream fusion and task-specific learning. The implemented strategy is two-stage and supports stage-specific checkpointing.

\subsection{Stage 1: Node Attribute Masking}
A fraction $\rho_{\text{node}}$ of nodes is masked in feature space. Let $\mathcal{M}$ denote masked node indices and $c_i$ the original atom-type class for node $i$. The stage-1 objective is:
\begin{equation}
\mathcal{L}_{\text{stage1}} = \frac{1}{|\mathcal{M}|} \sum_{i \in \mathcal{M}} \mathrm{CE}(\hat{\mathbf{p}}_i, c_i),
\end{equation}
where $\hat{\mathbf{p}}_i$ is the predicted atom-type distribution.

\subsection{Stage 2: Descriptor Prediction with Dynamic Masking}
The stage-2 model predicts graph-level descriptor targets. A dynamic supervision mask $\mathbf{M}$ is sampled per batch to keep only a fraction $\rho_{\text{keep}}$ of valid descriptor entries. If $\hat{\mathbf{z}}$ are predicted descriptors and $\mathbf{z}$ are targets:
\begin{equation}
\mathcal{L}_{\text{stage2}} = \frac{\sum_j M_j (\hat{z}_j - z_j)^2}{\sum_j M_j + \epsilon}.
\end{equation}

In addition, the implementation supports an optional graph-level contrastive objective (NT-Xent) between two stochastic embedding views. Let this term be denoted by $\mathcal{L}_{\text{graph}}$. The effective training objective is:
\begin{equation}
\mathcal{L}_{\text{total}} = \mathcal{L}_{\text{stage2}} + \lambda_{\text{graph}}\,\mathcal{L}_{\text{graph}},
\end{equation}
where $\lambda_{\text{graph}} \ge 0$ is the graph-objective weight. For $\lambda_{\text{graph}}=0$, training reduces to the masked descriptor objective alone.

This has two effects:
\begin{enumerate}[label=\alph*)]
    \item prevents over-reliance on fixed descriptor subsets,
    \item encourages robust latent chemistry features.
\end{enumerate}

\subsection{Target Preprocessing for Pretraining}
Descriptor targets are selected and standardized with robust safeguards:
\begin{itemize}
    \item remove unstable columns by minimum valid fraction and non-trivial variance,
    \item z-score normalization using finite entries,
    \item winsorization (clipping) to limit outlier impact,
    \item explicit valid-mask persistence for masked losses.
\end{itemize}

\subsection{Transfer into Downstream Models}
The encoder state is exported and reused downstream (e.g., TAME Stage-1 initialization). Stage labels and run metadata are stored in checkpoints to ensure consistent loading and avoid accidental stage mismatch.

\section{SEG: Structured Embedding Fusion (Graph + Text)}
\subsection{Architecture}
SEG combines:
\begin{enumerate}[label=\arabic*.]
    \item ChebNet graph encoder + graph pooling,
    \item optional text projection $\mathbb{R}^{d_t} \rightarrow \mathbb{R}^{d_t'}$,
    \item configurable fusion module,
    \item prediction head (linear or MLP).
\end{enumerate}

Given graph embedding $\mathbf{g}$ and text embedding $\mathbf{t}$, fusion produces
\begin{equation}
\mathbf{h}_{\text{fuse}} = \mathrm{Fuse}(\mathbf{g}, \mathbf{t}),
\end{equation}
followed by task head prediction.

\subsection{Fusion Variants}
The implementation supports multiple fusion families:
\begin{itemize}
    \item concat,
    \item cross-modal multi-head attention,
    \item gated fusion,
    \item FiLM-style conditioning,
    \item bilinear-style interaction.
\end{itemize}

In current benchmark use, cross-modal MHA is often selected as the primary SEG baseline.

\subsection{Loss Handling by Task Type}
\begin{itemize}
    \item \textbf{Single-task classification}: BCE.
    \item \textbf{Multi-label classification}: element-wise BCE with NaN masking.
    \item \textbf{Regression / multi-task regression}: MSE (masked when needed).
\end{itemize}

\section{TAME: Tri-Modal Element-Wise Mixture-of-Experts}
\subsection{Why TAME is the Largest Model}
TAME adds a third modality and introduces per-dimension tri-modal routing, increasing both representational capacity and adaptive behavior. Compared to SEG, it includes:
\begin{itemize}
    \item dedicated descriptor branch,
    \item router network over all three modalities,
    \item element-wise soft assignment per hidden dimension,
    \item optional balancing regularizers and modality dropout.
\end{itemize}

\subsection{Input Projections}
Each modality is projected to a shared hidden space $\mathbb{R}^{d_h}$:
\begin{align}
\mathbf{e}_g &= f_g(\mathbf{g}),\\
\mathbf{e}_t &= f_t(\mathbf{t}),\\
\mathbf{e}_d &= f_d(\mathbf{d}),
\end{align}
where $f_g, f_t$ are modality MLP projections and $f_d$ is descriptor projection with BatchNorm at input.

\subsection{Element-Wise Tri-Modal Routing}
Concatenate projections and route:
\begin{equation}
\mathbf{r} = R([\mathbf{e}_g;\mathbf{e}_t;\mathbf{e}_d]).
\end{equation}

Reshape router output to $(d_h, 3)$ per sample and apply softmax over modalities:
\begin{equation}
\mathbf{\Gamma}_{j,:} = \mathrm{softmax}(\mathbf{r}_{j,:}), \quad j=1,\dots,d_h,
\end{equation}
with gates $\gamma_{j,g}, \gamma_{j,t}, \gamma_{j,d}$.

Fused representation is computed element-wise:
\begin{equation}
\mathbf{h}_{\text{fuse},j} = \gamma_{j,g}\,\mathbf{e}_{g,j} + \gamma_{j,t}\,\mathbf{e}_{t,j} + \gamma_{j,d}\,\mathbf{e}_{d,j}.
\end{equation}

This is materially different from scalar modality gating: every hidden coordinate may favor a different modality mixture.

\subsection{Prediction Head and Task Generality}
TAME supports output dimensionality equal to number of tasks:
\begin{equation}
\hat{\mathbf{y}} = \mathrm{Head}(\mathbf{h}_{\text{fuse}}), \quad \hat{\mathbf{y}} \in \mathbb{R}^{T}.
\end{equation}

Hence, TAME handles:
\begin{itemize}
    \item $T=1$ single-task settings,
    \item $T>1$ multi-label classification,
    \item $T>1$ multi-task regression.
\end{itemize}

\subsection{Descriptor Branch Preprocessing}
For training descriptors, TAME performs:
\begin{enumerate}[label=\arabic*.]
    \item imputation by train-set medians,
    \item winsorization by configured quantiles $(q_{\ell}, q_u)$,
    \item optional standardization with train-set mean/std.
\end{enumerate}

This preprocessing is fit on training data only and reused for validation/test and inference.

\subsection{Regularization Terms Specific to TAME}
\paragraph{Descriptor modality dropout.}
During training, descriptor features can be randomly dropped (sample-wise) with probability $p_d$ to avoid over-dependence on descriptor branch.

\paragraph{Gate balance regularization.}
Let average gates per batch be $(\bar{g}, \bar{t}, \bar{d})$. A balancing penalty encourages non-collapsed usage:
\begin{equation}
\mathcal{L}_{\text{gate}} = (\bar{g}-\tfrac{1}{3})^2 + (\bar{t}-\tfrac{1}{3})^2 + (\bar{d}-\tfrac{1}{3})^2.
\end{equation}

Total loss:
\begin{equation}
\mathcal{L} = \mathcal{L}_{\text{task}} + \lambda_{\text{gate}}\,\mathcal{L}_{\text{gate}}.
\end{equation}

\section{Objectives by Regime}
\subsection{Classification}
Binary or multi-label classification uses BCE with logits. For missing labels (multi-label), let mask $M_{i,t} \in \{0,1\}$ indicate observed targets:
\begin{equation}
\mathcal{L}_{\text{BCE-mask}} = \frac{\sum_{i,t} M_{i,t}\,\ell_{\text{BCE}}(\hat{y}_{i,t}, y_{i,t})}{\sum_{i,t} M_{i,t} + \epsilon}.
\end{equation}

\subsection{Regression}
For multi-task regression:
\begin{equation}
\mathcal{L}_{\text{MSE-mask}} = \frac{\sum_{i,t} M_{i,t}(\hat{y}_{i,t}-y_{i,t})^2}{\sum_{i,t} M_{i,t} + \epsilon}.
\end{equation}
When all entries are observed, $M$ is all ones.

\section{Checkpointing and Transfer Safety}
The stack stores model state and metadata including stage labels and architecture settings. In practice, transfer routines enforce:
\begin{itemize}
    \item stage compatibility checks (e.g., stage1 encoder checkpoints for stage1 init paths),
    \item architecture consistency where strict loading is required,
    \item alias plus config-aware checkpoint discovery.
\end{itemize}

This reduces silent mismatch errors in long benchmark workflows.

\section{Evaluation Protocol Notes}
\subsection{Classification Benchmarks}
For BACE/Tox21-style classification, thresholds can be tuned on validation predictions and then applied to test predictions. Metrics commonly tracked:
\begin{itemize}
    \item ROC-AUC,
    \item Accuracy,
    \item F1,
    \item macro variants for multi-label settings.
\end{itemize}

\subsection{Regression Benchmarks}
For QM8-style multi-task regression:
\begin{itemize}
    \item train on normalized targets for optimization stability,
    \item inverse-transform predictions for final MAE/RMSE reporting,
    \item report both macro and per-task values.
\end{itemize}

\section{Practical Interpretation of Gate Statistics}
Reported gate means $(G/T/D)$ per epoch provide a coarse but useful diagnostic:
\begin{itemize}
    \item near-uniform usage often indicates balanced modality contribution,
    \item severe collapse to one branch may indicate overfitting or modality dominance,
    \item stage1-initialized graph encoders can shift gate allocation toward graph branch while preserving descriptor utility.
\end{itemize}

Because routing is element-wise, mean gates summarize trends but do not capture full per-dimension specialization.

\section{Design Strengths and Known Tradeoffs}
\subsection{Strengths}
\begin{itemize}
    \item Unified treatment of single-task and multi-task settings.
    \item Explicit missing-label masking for robust multi-label training.
    \item Descriptor branch and element-wise routing improve representational richness.
    \item Stage-aware pretraining and transfer path improves engineering reliability.
\end{itemize}

\subsection{Tradeoffs}
\begin{itemize}
    \item TAME has higher complexity and tuning burden than SEG.
    \item Task gains from stage1 initialization can be dataset-dependent.
    \item Macro metrics can hide target-specific variance; per-task diagnostics remain essential.
\end{itemize}

\section{Recommended Reporting Template for Publications}
For a high-quality scientific report, include:
\begin{enumerate}[label=\arabic*.]
    \item Architecture diagram with modality paths and router shape.
    \item Full hyperparameter tables for SEG and TAME.
    \item Pretraining details (stage objectives, masking rates, data scale).
    \item Validation protocol (early stopping, threshold policy).
    \item Main table: macro metrics plus confidence intervals across seeds.
    \item Appendix: per-task metrics, gate trajectories, and ablation studies.
\end{enumerate}

\section{Suggested Ablations}
\begin{itemize}
    \item Remove descriptor branch (TAME to SEG-like behavior).
    \item Replace element-wise router with scalar modality gates.
    \item Stage1 init on/off under identical tuning.
    \item Gate-balance weight sweep.
    \item Descriptor dropout sweep.
    \item Pooling choice sensitivity (mean vs set2set).
\end{itemize}

\section{Conclusion}
The architecture family forms a coherent progression:
\begin{itemize}
    \item \textbf{ChebNet pretraining} provides transferable chemical graph priors.
    \item \textbf{SEG} supplies a strong bi-modal baseline with modular fusion.
    \item \textbf{TAME} extends this to tri-modal, element-wise routed fusion and currently represents the most expressive model in the stack.
\end{itemize}

In benchmark practice, TAME is the primary high-capacity architecture, while SEG remains an efficient and informative baseline. Stage-aware ChebNet pretraining contributes a principled initialization pathway and an interpretable bridge between unsupervised chemistry pretraining and supervised multimodal prediction.

\appendix
\section{Minimal Reproducibility Checklist}
\begin{itemize}
    \item Fix global seed for NumPy and PyTorch.
    \item Log exact split type and split seed.
    \item Log checkpoint path and pretraining stage label.
    \item Store full hyperparameter dictionaries for each run.
    \item Report metrics with macro and per-task breakdowns where applicable.
\end{itemize}

\end{document}
